%------------------------------------------------------------------------------%
% Luis A. D'Afonseca
%
%------------------------------------------------------------------------------%
\documentclass[a4paper,12pt,fleqn]{article}

\usepackage{style_avaliacao}

\ShowAnswers

%------------------------------------------------------------------------------%
\begin{document}

\makeHeader{Integrais e Séries}{Prova 3 -- Turma 4}{09/11/2023}

% 01
%------------------------------------------------------------------------------%
\exercise{25}
Use o teste da integral para verificar se a soma da série existe
\[
  \sum_{n=1}^{\infty} \frac{5}{\left(1+n\right)^2}
\]
Não se esqueça de verificar que as condições do teste são satisfeitas.

\begin{answer}
  Vamos usar a função real
  \[
    f(x) = \frac{5}{\left(1+x\right)^2}
  \]

  Precisamos verificar as três condições do teste
  \begin{enumerate}[label=\roman*)]
    \item \(f(x)>0\) em \([1, \infty)\) \\[2mm]
      Como \((1+x)^2>0\) para todo \(x\in\R\) então \(f(x)>0\)

    \item \(f(x)\) é contínua em \([1, \infty)\) \\[2mm]
      Como \((1+x)^2>0\) a função \(f(x)\) é contínua em \(\R\)

    \item \(f(x)\) é decrescente em  \([1, \infty)\) \\[2mm]
      A derivada de \(f(x)\) é
      \begin{align*}
        f'(x)
        & = \frac
          {(5)'\left(1+x\right)^2 - 5\left[\left(1+x\right)^2\right]'}
          {\left(1+x\right)^4} \\[2mm]
        & = \frac{-5\times 2\left(1+x\right) 1}{\left(1+x\right)^4} \\[2mm]
        & = \frac{-10\left(1+x\right)}{\left(1+x\right)^4} \\[2mm]
        & = \frac{-10}{\left(1+x\right)^3}
      \end{align*}
      portanto \(f'(x)<0\) para \(x>-1\), assim \(f\) é decrescente no intervalo \([1, \infty)\)
  \end{enumerate}
  Calculando a primitiva de \(f\)
  \[
    F(x) = \int \frac{5}{\left(1+x\right)^2} dx
  \]
  Fazendo a mudança de variáveis
  \[
    u = 1+x
    \qquad
    du = dx
  \]
  \[
    F(x)
     = \int \frac{5}{\left(1+x\right)^2} dx
     = 5 \int \frac{du}{u^2}
     = 5 \int u^{-2} du
     = 5 \frac{u^{-1}}{-1} + c
     = \frac{-5}{u} + c
     = \frac{-5}{1+x} + c
  \]
  Calculando a integral imprópria
  \[
    A
    = \int_1^\infty f(x)dx
    = \lim_{b\to\infty}\int_1^b f(x)dx
    = \lim_{b\to\infty} F(x)\evalat{1}{b}
    = \lim_{b\to\infty} \left(\frac{-5}{1+b} - \frac{-5}{1+1}\right)
    = \frac{5}{2}
  \]
  Como a integral converge a série também converge e portanto sua soma existe.
\end{answer}

% 02
%------------------------------------------------------------------------------%
\exercise{25}
Calcule o limite de \(\left(\frac{\cos(n) + 1}{\ln(n)}\right)_{n=2}^\infty\)

\begin{answer}
  Queremos calcular
  \[
    \lim_{n\to\infty}\frac{\cos(n) + 1}{\ln(n)}
  \]
  Como \(\lim_{n\to\infty}\cos(n)\) não existe,
  não podemos calcular o limite de cada termo
  separadamente, portanto, vamos usar o Teorema do Confronto.
  \begin{align*}
    -1 \leq & \cos(n) \leq 1 \\[2mm]
     0 \leq & \cos(n)+1 \leq 2 \\[2mm]
     \frac{0}{\,\ln(n)\,} \leq & \frac{\,\cos(n)+1\,}{\ln(n)} \leq \frac{2}{\,\ln(n)\,} \\[2mm]
     0 \leq & \frac{\,\cos(n)+1\,}{\ln(n)} \leq \frac{2}{\,\ln(n)\,}
  \end{align*}
  A penúltima desigualdade vale, pois, \(\ln(n)>0\) para todo \(n\geq2\).

  Como \(\lim_{n\to\infty} 0 = \lim_{n\to\infty} \frac{2}{\ln(n)}=0\),
  concluímos que
  \[
    \lim_{n\to\infty}\frac{\cos(n) + 1}{\ln(n)} = 0
  \]
\end{answer}

% 03
%------------------------------------------------------------------------------%
\exercise{25}
Verifique se a série converge
\[
  \sum_{n=1}^\infty \frac{3^n}{3+9^n}
\]

\begin{answer}
  Como todos os termos são positivos, podemos usar o teste da comparação no limite.

  Escolhemos \(b_n=\sfrac{1}{3^n}\), que são os termos de uma série
  geométrica convergente, pois \(\abs{r}=\sfrac{1}{3}<1\).

  \begin{align*}
    c
    & = \lim_{n\to\infty} \frac{a_n}{b_n} \\[2mm]
    & = \lim_{n\to\infty} \frac{\;\dfrac{3^n}{3+9^n}\;}{\;\dfrac{1}{3^n}\;} \\[2mm]
    & = \lim_{n\to\infty} \frac{3^n}{3+9^n}\;3^n \\[2mm]
    & = \lim_{n\to\infty} \frac{9^n}{3+9^n} \\[2mm]
    & = \lim_{n\to\infty} \frac{1}{\sfrac{3}{9^n}+1} \\[2mm]
    & = 1
  \end{align*}
  Como \(c>0\) e finito, o teste da comparação no limite diz que ou ambas as séries
  convergem ou ambas divergem. Como a soma de \(b_n\) converge, então
  \[
    \sum_{n=1}^\infty \frac{3^n}{3+9^n}
  \]
  também converge.
\end{answer}

% 04
%------------------------------------------------------------------------------%
\exercise{25}
Considerando a série
\[
  \sum_{n=1}^{\infty} \frac{(-2)^n}{3^{n-1}}
\]
\begin{enumerate}[label=\alph*)]
  \item Prove que a série converge
  \item Calcule a soma da série
  \item Estime o erro cometido ao aproximar a soma por \(S_{100}\)
\end{enumerate}

\begin{answer}
\begin{enumerate}[label=\alph*)]
  \item Essa é uma série geométrica com
    \(\alpha = -2\) e
    \(r=\frac{-2}{3}\),
    pois,
    \[
      \sum_{n=1}^{\infty} \frac{(-2)^n}{3^{n-1}}
      = \sum_{n=1}^{\infty} \frac{(-2)(-2)^{n-1}}{3^{n-1}}
      = \sum_{n=1}^{\infty} (-2)\left(\frac{-2}{3}\right)^{n-1}
    \]
    Como
    \(\abs{r} = \frac{2}{3} < 1\)
    a série é convergente.
  \item
    A soma da série geométrica é
    \[
      S
      = \frac{\alpha}{1-r}
      = \frac{-2}{1-\left(\frac{-2}{3}\right)}
      = \frac{-2}{\frac{3}{3}+\frac{2}{3}}
      = \frac{-2}{\frac{5}{3}}
      = -2\frac{3}{5}
      = \frac{-6}{5}
    \]
  \item Além de ser uma série geométrica essa série também é alternada, então
  \[
    R_{100} < \abs{a_{101}}
  \]
  como \(\abs{a_n}=\frac{2^n}{3^{n-1}}\), portanto
  \[
    R_{100} < \frac{2^{101}}{3^{100}}
    \approx 4,9193 \times 10^{-18}
  \]
\end{enumerate}
\end{answer}

%------------------------------------------------------------------------------%
\end{document}
%------------------------------------------------------------------------------%
